% Generated. Do not edit.
% Projection scriptural-date-location/section, version 1: the Scriptural Date and
% Location page, generated whole from research/chronology.toml, the Scripture
% chronology corpus and the source library by tools/tpt document-types refresh.
% document: liturgy/roman-rite/postconciliar/roman-missal-third-edition-en-us-2011/propers/temporal/pc-s53-twenty-seventh-sunday-in-ordinary-time-year-a
\runninghead{Scriptural Date and Location}
\section{Scriptural Date and Location}
\label{sec:date-location}

{\small Each passage the Mass reads or sings from Scripture, in canonical order. The Context column names what is dated --- the event the passage narrates, a setting it remembers, the traditional attribution of its author, or the composition of its book --- and where it belongs; the Date column gives the answer of the Scripture chronology corpus for that subject, with every alternative it records. Texts of the Mass that are not drawn from Scripture have no row. The evidence for each date follows the table.\par}
\medskip
% Generated. Do not edit.
% Deterministic projection of research/chronology.toml.
% schema: 5
% document: liturgy/roman-rite/postconciliar/roman-missal-third-edition-en-us-2011/propers/temporal/pc-s53-twenty-seventh-sunday-in-ordinary-time-year-a
% generated-by: tools/tpt proper-chronology annotations
\newcommand{\chronologyannotationclaim}[6]{#6}
\newcommand{\chronologyannotationreach}[3]{}
\newcommand{\chronologyannotationgroup}[3]{#3}
\newcommand{\chronologyannotation}[1]{%
  \ifcsname triptychchronologyannotation@#1\endcsname
    \csname triptychchronologyannotation@#1\endcsname
  \else
    \PackageError{triptych}{No chronology annotation for #1}{}%
  \fi}
% appointment-dependency: guidance/liturgy/postconciliar-propers-registry.md sha256=376860a7c37a4c75bb6d880cd1b260183be1877a07e3cd01660a9e2854e2ea72
% appointment-dependency: src/gpt/liturgy/roman-rite/postconciliar/roman-missal-third-edition-en-us-2011/propers/registry/README.md sha256=4e2d92326dece103e1475e70ffb5e447815f9b12d697926e48a54850a1641208
% appointment-dependency: src/gpt/liturgy/roman-rite/postconciliar/roman-missal-third-edition-en-us-2011/propers/registry/formula-dispositions.md sha256=de58754f5a8ab61dffa3f4a6a479c843fe3763b6a5b37099282960de71d93989
% appointment-dependency: src/gpt/liturgy/roman-rite/postconciliar/roman-missal-third-edition-en-us-2011/propers/temporal/pc-s53-twenty-seventh-sunday-in-ordinary-time-year-a/instance/manifest.md sha256=9ab2f377d6de66eb639a2115620d32379b1d22b40cea1ac67dcf429b821630e9
% appointment-dependency: src/gpt/liturgy/roman-rite/postconciliar/roman-missal-third-edition-en-us-2011/propers/temporal/pc-s53-twenty-seventh-sunday-in-ordinary-time-year-a/propers/verified.md sha256=c4367c88e48ca6954481a904edde9c857128f2a94e74649cc38e993cafb22bb8
% appointment-dependency: src/gpt/liturgy/roman-rite/postconciliar/roman-missal-third-edition-en-us-2011/propers/temporal/pc-s53-twenty-seventh-sunday-in-ordinary-time-year-a/research/chronology-inputs.toml sha256=9fd487f1561888df70c71ff005090c7a9ff3a09bf7d846ab70671ac156997841
% appointment-dependency: src/gpt/liturgy/roman-rite/postconciliar/roman-missal-third-edition-en-us-2011/propers/temporal/shared/ordinary-time/weeks/27/propers/verified.md sha256=49ce6830cc5bce8d654908675720c69c73ad0e44ee3832e3196288406a137cd6
% appointment-note: entrance: Esther 13:9-11; identified-basis; vulgate numbering; canonical verse queries, not appointed wording.
% appointment-note: collect: Composed Missal oration; no directly appointed biblical citation.
% appointment-note: first-reading: Isaiah 5:1-7; appointed; vulgate numbering; canonical verse queries, not appointed wording.
% appointment-note: responsorial-psalm: Psalm 80:9, 12, 13-14, 15-16, 19-20; appointed; hebrew numbering; canonical verse queries, not appointed wording.
% appointment-note: responsorial-psalm: Isaiah 5:7a; response; vulgate numbering; whole-verse query envelope, not subverse chronology or appointed wording.
% appointment-note: second-reading: Philippians 4:6-9; appointed; vulgate numbering; canonical verse queries, not appointed wording.
% appointment-note: acclamation: Cf. John 15:16; adaptation; vulgate numbering; canonical verse queries, not appointed wording.
% appointment-note: gospel: Matthew 21:33-43; appointed; vulgate numbering; canonical verse queries, not appointed wording.
% appointment-note: offerings: Composed Missal oration; no directly appointed biblical citation.
% appointment-note: communion-lamentations: Lamentations 3:25; appointed; vulgate numbering; canonical verse queries, not appointed wording.
% appointment-note: communion-corinthians: Cf. 1 Corinthians 10:17; adaptation; vulgate numbering; canonical verse queries, not appointed wording.
% appointment-note: communion-corinthians: 1 Corinthians 10:16; identified-basis; vulgate numbering; canonical verse queries, not appointed wording.
% appointment-note: after-communion: Composed Missal oration; no directly appointed biblical citation.
% entrance: Entrance antiphon
\expandafter\def\csname triptychchronologyannotation@entrance\endcsname{%
  \chronologyannotationgroup{composition}{preferred}{\textbf{Composition}: \chronologyannotationreach{Esth.13.10}{Esth}{true}\chronologyannotationreach{Esth.13.11}{Esth}{true}\chronologyannotationreach{Esth.13.9}{Esth}{true}\chronologyannotationclaim{composition.book-of-esther}{composition}{catholic-traditional-v1}{preferred}{at the end of the reign of Xerxes I (485-465 B.C.) or during the reign of his son Artaxerxes I (465-425 B.C.)}{B.C. 465--425}.}%
}
% first-reading: First reading
\expandafter\def\csname triptychchronologyannotation@first-reading\endcsname{%
  \chronologyannotationgroup{traditional-attribution}{preferred}{\textbf{Traditional attribution}: \chronologyannotationreach{Is.5.1}{Is}{true}\chronologyannotationreach{Is.5.2}{Is}{true}\chronologyannotationreach{Is.5.3}{Is}{true}\chronologyannotationreach{Is.5.4}{Is}{true}\chronologyannotationreach{Is.5.5}{Is}{true}\chronologyannotationreach{Is.5.6}{Is}{true}\chronologyannotationreach{Is.5.7}{Is}{true}\chronologyannotationclaim{israel.prophets.isaias-traditional-era}{traditional-attribution}{catholic-traditional-v1}{preferred}{from the closing year of Ozias, King of Juda (740 B.C.); no prophecies appear to be later than 701 B.C.}{Isaias (ministry in Souvay's traditional account), B.C. 740--701}.}%
}
% responsorial-psalm: Responsorial psalm
\expandafter\def\csname triptychchronologyannotation@responsorial-psalm\endcsname{%
  \chronologyannotationgroup{historical-setting}{research-pending}{\textbf{Historical setting} -- Occasion date unresolved.}%
  \space \chronologyannotationgroup{composition}{nonuniform}{\textbf{Composition} -- No single date applies across every cited locus.}%
}
% second-reading: Second reading
\expandafter\def\csname triptychchronologyannotation@second-reading\endcsname{%
  \chronologyannotationgroup{composition}{disputed}{\textbf{Composition} -- disputed: \chronologyannotationreach{Phil.4.6}{Phil}{true}\chronologyannotationreach{Phil.4.7}{Phil}{true}\chronologyannotationreach{Phil.4.8}{Phil}{true}\chronologyannotationreach{Phil.4.9}{Phil}{true}\chronologyannotationclaim{composition.epistle-to-the-philippians}{composition}{catholic-traditional-v1}{disputed}{(Philemon; Colossians; Ephesians; Philippians), 61}{A.D. 61}; \chronologyannotationreach{Phil.4.6}{Phil}{true}\chronologyannotationreach{Phil.4.7}{Phil}{true}\chronologyannotationreach{Phil.4.8}{Phil}{true}\chronologyannotationreach{Phil.4.9}{Phil}{true}\chronologyannotationclaim{composition.epistle-to-the-philippians}{composition}{catholic-traditional-v1}{disputed}{at Rome (A.D. 62-64)}{A.D. 62--64}.}%
}
% acclamation: Gospel acclamation
\expandafter\def\csname triptychchronologyannotation@acclamation\endcsname{%
  \chronologyannotationgroup{narrated-event}{disputed}{\textbf{Event} -- disputed: \chronologyannotationreach{John.15.16}{Matt.26.20-29 Mark.14.17-25 Luke.22.14-38 John.13.1- John.14 John.15 John.16 John.17.1-26}{true}\chronologyannotationclaim{life-of-christ.last-supper}{narrated-event}{catholic-traditional-v1}{disputed}{the fourteenth day of Nisan}{The Last Supper, 01-14}.}%
  \space \chronologyannotationgroup{composition}{disputed}{\textbf{Composition} -- disputed: \chronologyannotationreach{John.15.16}{John}{true}\chronologyannotationclaim{composition.gospel-of-john}{composition}{catholic-traditional-v1}{disputed}{from the year 90 to 100 (approximately)}{c. A.D. 90--100}; \chronologyannotationreach{John.15.16}{John}{true}\chronologyannotationclaim{composition.gospel-of-john}{composition}{catholic-traditional-v1}{disputed}{96 or one of the succeeding years}{A.D. 96--100}.}%
}
% gospel: Gospel
\expandafter\def\csname triptychchronologyannotation@gospel\endcsname{%
  \chronologyannotationgroup{narrated-event}{research-pending}{\textbf{Event} -- No narrated-event date in the chronology corpus.}%
  \space \chronologyannotationgroup{composition}{disputed}{\textbf{Composition} -- disputed: \chronologyannotationreach{Matt.21.33}{Matt}{true}\chronologyannotationreach{Matt.21.34}{Matt}{true}\chronologyannotationreach{Matt.21.35}{Matt}{true}\chronologyannotationreach{Matt.21.36}{Matt}{true}\chronologyannotationreach{Matt.21.37}{Matt}{true}\chronologyannotationreach{Matt.21.38}{Matt}{true}\chronologyannotationreach{Matt.21.39}{Matt}{true}\chronologyannotationreach{Matt.21.40}{Matt}{true}\chronologyannotationreach{Matt.21.41}{Matt}{true}\chronologyannotationreach{Matt.21.42}{Matt}{true}\chronologyannotationreach{Matt.21.43}{Matt}{true}\chronologyannotationclaim{composition.gospel-of-matthew}{composition}{catholic-traditional-v1}{disputed}{about A.D. 38-45}{c. A.D. 38--45}; \chronologyannotationreach{Matt.21.33}{Matt}{true}\chronologyannotationreach{Matt.21.34}{Matt}{true}\chronologyannotationreach{Matt.21.35}{Matt}{true}\chronologyannotationreach{Matt.21.36}{Matt}{true}\chronologyannotationreach{Matt.21.37}{Matt}{true}\chronologyannotationreach{Matt.21.38}{Matt}{true}\chronologyannotationreach{Matt.21.39}{Matt}{true}\chronologyannotationreach{Matt.21.40}{Matt}{true}\chronologyannotationreach{Matt.21.41}{Matt}{true}\chronologyannotationreach{Matt.21.42}{Matt}{true}\chronologyannotationreach{Matt.21.43}{Matt}{true}\chronologyannotationclaim{composition.gospel-of-matthew}{composition}{catholic-traditional-v1}{disputed}{about the year 40-42}{c. A.D. 40--42}; \chronologyannotationreach{Matt.21.33}{Matt}{true}\chronologyannotationreach{Matt.21.34}{Matt}{true}\chronologyannotationreach{Matt.21.35}{Matt}{true}\chronologyannotationreach{Matt.21.36}{Matt}{true}\chronologyannotationreach{Matt.21.37}{Matt}{true}\chronologyannotationreach{Matt.21.38}{Matt}{true}\chronologyannotationreach{Matt.21.39}{Matt}{true}\chronologyannotationreach{Matt.21.40}{Matt}{true}\chronologyannotationreach{Matt.21.41}{Matt}{true}\chronologyannotationreach{Matt.21.42}{Matt}{true}\chronologyannotationreach{Matt.21.43}{Matt}{true}\chronologyannotationclaim{composition.gospel-of-matthew}{composition}{catholic-traditional-v1}{disputed}{the years 40-45}{A.D. 40--45}; \chronologyannotationreach{Matt.21.33}{Matt}{true}\chronologyannotationreach{Matt.21.34}{Matt}{true}\chronologyannotationreach{Matt.21.35}{Matt}{true}\chronologyannotationreach{Matt.21.36}{Matt}{true}\chronologyannotationreach{Matt.21.37}{Matt}{true}\chronologyannotationreach{Matt.21.38}{Matt}{true}\chronologyannotationreach{Matt.21.39}{Matt}{true}\chronologyannotationreach{Matt.21.40}{Matt}{true}\chronologyannotationreach{Matt.21.41}{Matt}{true}\chronologyannotationreach{Matt.21.42}{Matt}{true}\chronologyannotationreach{Matt.21.43}{Matt}{true}\chronologyannotationclaim{composition.gospel-of-matthew}{composition}{catholic-traditional-v1}{disputed}{about the year 60-68}{c. A.D. 60--68}; \chronologyannotationreach{Matt.21.33}{Matt}{true}\chronologyannotationreach{Matt.21.34}{Matt}{true}\chronologyannotationreach{Matt.21.35}{Matt}{true}\chronologyannotationreach{Matt.21.36}{Matt}{true}\chronologyannotationreach{Matt.21.37}{Matt}{true}\chronologyannotationreach{Matt.21.38}{Matt}{true}\chronologyannotationreach{Matt.21.39}{Matt}{true}\chronologyannotationreach{Matt.21.40}{Matt}{true}\chronologyannotationreach{Matt.21.41}{Matt}{true}\chronologyannotationreach{Matt.21.42}{Matt}{true}\chronologyannotationreach{Matt.21.43}{Matt}{true}\chronologyannotationclaim{composition.gospel-of-matthew}{composition}{catholic-traditional-v1}{disputed}{about the years 64-67}{c. A.D. 64--67}; \chronologyannotationreach{Matt.21.33}{Matt}{true}\chronologyannotationreach{Matt.21.34}{Matt}{true}\chronologyannotationreach{Matt.21.35}{Matt}{true}\chronologyannotationreach{Matt.21.36}{Matt}{true}\chronologyannotationreach{Matt.21.37}{Matt}{true}\chronologyannotationreach{Matt.21.38}{Matt}{true}\chronologyannotationreach{Matt.21.39}{Matt}{true}\chronologyannotationreach{Matt.21.40}{Matt}{true}\chronologyannotationreach{Matt.21.41}{Matt}{true}\chronologyannotationreach{Matt.21.42}{Matt}{true}\chronologyannotationreach{Matt.21.43}{Matt}{true}\chronologyannotationclaim{composition.gospel-of-matthew}{composition}{catholic-traditional-v1}{disputed}{about the year 50}{c. A.D. 50}.}%
}
% communion-lamentations: Communion antiphon: Lamentations option
\expandafter\def\csname triptychchronologyannotation@communion-lamentations\endcsname{%
  \chronologyannotationgroup{composition}{preferred}{\textbf{Composition}: \chronologyannotationreach{Lam.3.25}{Lam}{true}\chronologyannotationclaim{critical.lamentations.after-jerusalem-destruction}{composition}{catholic-critical-v1}{preferred}{after Jerusalem's destruction in 587 B.C.; grouped later at an unknown date}{Responses to Jerusalem's destruction in 587 B.C.}}%
}
% communion-corinthians: Communion antiphon: Corinthians option
\expandafter\def\csname triptychchronologyannotation@communion-corinthians\endcsname{%
  \chronologyannotationgroup{composition}{disputed}{\textbf{Composition} -- disputed: \chronologyannotationreach{1Cor.10.16}{1Cor}{true}\chronologyannotationreach{1Cor.10.17}{1Cor}{true}\chronologyannotationclaim{composition.first-epistle-to-the-corinthians}{composition}{catholic-traditional-v1}{disputed}{(1 and 2 Corinthians; Galatians), 56}{A.D. 56}; \chronologyannotationreach{1Cor.10.16}{1Cor}{true}\chronologyannotationreach{1Cor.10.17}{1Cor}{true}\chronologyannotationclaim{composition.first-epistle-to-the-corinthians}{composition}{catholic-traditional-v1}{disputed}{about Easter A.D. 57}{A.D. 57}.}%
}
% presentation: orientation-v3
\newcommand{\chronologyorientationrow}[2]{%
  \ifcsname triptychchronologyorientation@#1\endcsname
    \csname triptychchronologyorientation@#1\endcsname{#2}%
  \else
    \PackageError{triptych}{No chronology orientation row for #1}{}%
  \fi}
\expandafter\def\csname triptychchronologyorientation@row-01242fba70f5222d\endcsname#1{%
Entrance antiphon: Esth. 13:9--11 & Composition: Esther; story at Susa; writing in Persia conjectured & \chronologyannotationreach{Esth.13.10}{Esth}{true}\chronologyannotationreach{Esth.13.11}{Esth}{true}\chronologyannotationreach{Esth.13.9}{Esth}{true}\chronologyannotationclaim{composition.book-of-esther}{composition}{catholic-traditional-v1}{preferred}{at the end of the reign of Xerxes I (485-465 B.C.) or during the reign of his son Artaxerxes I (465-425 B.C.)}{B.C.~465--425 (Conjectural early dating; later date possible)}\\[.3em]%
}
\expandafter\def\csname triptychchronologyorientation@row-2a5f76fbcf6ad0bd\endcsname#1{%
Responsorial psalm: Ps. 79:9,12--16,19--20 (Vg.; Heb. 80) & Historical setting & Occasion date unresolved\\*[.1em]%
}
\expandafter\def\csname triptychchronologyorientation@row-af9138e4ceea5e79\endcsname#1{%
 & Composition: Psalms & \chronologyannotationreach{Ps.79.12}{Ps}{true}\chronologyannotationreach{Ps.79.13}{Ps}{true}\chronologyannotationreach{Ps.79.14}{Ps}{true}\chronologyannotationreach{Ps.79.15}{Ps}{true}\chronologyannotationreach{Ps.79.16}{Ps}{true}\chronologyannotationreach{Ps.79.19}{Ps}{true}\chronologyannotationreach{Ps.79.20}{Ps}{true}\chronologyannotationreach{Ps.79.9}{Ps}{true}\chronologyannotationclaim{critical.psalms.latest-composition-boundary}{composition}{catholic-critical-v1}{preferred}{before the Maccabean period, around 165 B.C.; no individual psalm can be dated securely}{Before c.~165 B.C.}\\[.3em]%
}
\expandafter\def\csname triptychchronologyorientation@row-34474a8e5d3d77f6\endcsname#1{%
First reading: Is. 5:1--7 & Traditional attribution: Isaias (ministry in Souvay's traditional account); Judah and Jerusalem & \chronologyannotationreach{Is.5.1}{Is}{true}\chronologyannotationreach{Is.5.2}{Is}{true}\chronologyannotationreach{Is.5.3}{Is}{true}\chronologyannotationreach{Is.5.4}{Is}{true}\chronologyannotationreach{Is.5.5}{Is}{true}\chronologyannotationreach{Is.5.6}{Is}{true}\chronologyannotationreach{Is.5.7}{Is}{true}\chronologyannotationclaim{israel.prophets.isaias-traditional-era}{traditional-attribution}{catholic-traditional-v1}{preferred}{from the closing year of Ozias, King of Juda (740 B.C.); no prophecies appear to be later than 701 B.C.}{B.C.~740--701}\\[.3em]%
}
\expandafter\def\csname triptychchronologyorientation@row-58ef6f1ff512a823\endcsname#1{%
Responsorial psalm: Is. 5:7 & Traditional attribution: Isaias (ministry in Souvay's traditional account); Judah and Jerusalem & \chronologyannotationreach{Is.5.7}{Is}{true}\chronologyannotationclaim{israel.prophets.isaias-traditional-era}{traditional-attribution}{catholic-traditional-v1}{preferred}{from the closing year of Ozias, King of Juda (740 B.C.); no prophecies appear to be later than 701 B.C.}{B.C.~740--701}\\[.3em]%
}
\expandafter\def\csname triptychchronologyorientation@row-8dc930e93054888b\endcsname#1{%
Communion antiphon: Lamentations option: Lam. 3:25 & Composition: The poems of Lamentations after Jerusalem's destruction & \chronologyannotationreach{Lam.3.25}{Lam}{true}\chronologyannotationclaim{critical.lamentations.after-jerusalem-destruction}{composition}{catholic-critical-v1}{preferred}{after Jerusalem's destruction in 587 B.C.; grouped later at an unknown date}{Responses to Jerusalem's destruction in 587 B.C.}\\[.3em]%
}
\expandafter\def\csname triptychchronologyorientation@row-a6eb6120b16899b0\endcsname#1{%
Gospel: Matt. 21:33--43 & Narrated event & Date unresolved\\*[.1em]%
}
\expandafter\def\csname triptychchronologyorientation@row-0bfa70ac74041a8c\endcsname#1{%
 & Composition: Matthew; Semitic original associated with Palestine & Disputed: \chronologyannotationreach{Matt.21.33}{Matt}{true}\chronologyannotationreach{Matt.21.34}{Matt}{true}\chronologyannotationreach{Matt.21.35}{Matt}{true}\chronologyannotationreach{Matt.21.36}{Matt}{true}\chronologyannotationreach{Matt.21.37}{Matt}{true}\chronologyannotationreach{Matt.21.38}{Matt}{true}\chronologyannotationreach{Matt.21.39}{Matt}{true}\chronologyannotationreach{Matt.21.40}{Matt}{true}\chronologyannotationreach{Matt.21.41}{Matt}{true}\chronologyannotationreach{Matt.21.42}{Matt}{true}\chronologyannotationreach{Matt.21.43}{Matt}{true}\chronologyannotationclaim{composition.gospel-of-matthew}{composition}{catholic-traditional-v1}{disputed}{about A.D. 38-45}{c.~A.D.~38--45}; \chronologyannotationreach{Matt.21.33}{Matt}{true}\chronologyannotationreach{Matt.21.34}{Matt}{true}\chronologyannotationreach{Matt.21.35}{Matt}{true}\chronologyannotationreach{Matt.21.36}{Matt}{true}\chronologyannotationreach{Matt.21.37}{Matt}{true}\chronologyannotationreach{Matt.21.38}{Matt}{true}\chronologyannotationreach{Matt.21.39}{Matt}{true}\chronologyannotationreach{Matt.21.40}{Matt}{true}\chronologyannotationreach{Matt.21.41}{Matt}{true}\chronologyannotationreach{Matt.21.42}{Matt}{true}\chronologyannotationreach{Matt.21.43}{Matt}{true}\chronologyannotationclaim{composition.gospel-of-matthew}{composition}{catholic-traditional-v1}{disputed}{about the year 40-42}{c.~A.D.~40--42}; \chronologyannotationreach{Matt.21.33}{Matt}{true}\chronologyannotationreach{Matt.21.34}{Matt}{true}\chronologyannotationreach{Matt.21.35}{Matt}{true}\chronologyannotationreach{Matt.21.36}{Matt}{true}\chronologyannotationreach{Matt.21.37}{Matt}{true}\chronologyannotationreach{Matt.21.38}{Matt}{true}\chronologyannotationreach{Matt.21.39}{Matt}{true}\chronologyannotationreach{Matt.21.40}{Matt}{true}\chronologyannotationreach{Matt.21.41}{Matt}{true}\chronologyannotationreach{Matt.21.42}{Matt}{true}\chronologyannotationreach{Matt.21.43}{Matt}{true}\chronologyannotationclaim{composition.gospel-of-matthew}{composition}{catholic-traditional-v1}{disputed}{the years 40-45}{A.D.~40--45}; \chronologyannotationreach{Matt.21.33}{Matt}{true}\chronologyannotationreach{Matt.21.34}{Matt}{true}\chronologyannotationreach{Matt.21.35}{Matt}{true}\chronologyannotationreach{Matt.21.36}{Matt}{true}\chronologyannotationreach{Matt.21.37}{Matt}{true}\chronologyannotationreach{Matt.21.38}{Matt}{true}\chronologyannotationreach{Matt.21.39}{Matt}{true}\chronologyannotationreach{Matt.21.40}{Matt}{true}\chronologyannotationreach{Matt.21.41}{Matt}{true}\chronologyannotationreach{Matt.21.42}{Matt}{true}\chronologyannotationreach{Matt.21.43}{Matt}{true}\chronologyannotationclaim{composition.gospel-of-matthew}{composition}{catholic-traditional-v1}{disputed}{about the year 60-68}{c.~A.D.~60--68}; \chronologyannotationreach{Matt.21.33}{Matt}{true}\chronologyannotationreach{Matt.21.34}{Matt}{true}\chronologyannotationreach{Matt.21.35}{Matt}{true}\chronologyannotationreach{Matt.21.36}{Matt}{true}\chronologyannotationreach{Matt.21.37}{Matt}{true}\chronologyannotationreach{Matt.21.38}{Matt}{true}\chronologyannotationreach{Matt.21.39}{Matt}{true}\chronologyannotationreach{Matt.21.40}{Matt}{true}\chronologyannotationreach{Matt.21.41}{Matt}{true}\chronologyannotationreach{Matt.21.42}{Matt}{true}\chronologyannotationreach{Matt.21.43}{Matt}{true}\chronologyannotationclaim{composition.gospel-of-matthew}{composition}{catholic-traditional-v1}{disputed}{about the years 64-67}{c.~A.D.~64--67}; \chronologyannotationreach{Matt.21.33}{Matt}{true}\chronologyannotationreach{Matt.21.34}{Matt}{true}\chronologyannotationreach{Matt.21.35}{Matt}{true}\chronologyannotationreach{Matt.21.36}{Matt}{true}\chronologyannotationreach{Matt.21.37}{Matt}{true}\chronologyannotationreach{Matt.21.38}{Matt}{true}\chronologyannotationreach{Matt.21.39}{Matt}{true}\chronologyannotationreach{Matt.21.40}{Matt}{true}\chronologyannotationreach{Matt.21.41}{Matt}{true}\chronologyannotationreach{Matt.21.42}{Matt}{true}\chronologyannotationreach{Matt.21.43}{Matt}{true}\chronologyannotationclaim{composition.gospel-of-matthew}{composition}{catholic-traditional-v1}{disputed}{about the year 50}{c.~A.D.~50}\\[.3em]%
}
\expandafter\def\csname triptychchronologyorientation@row-032e4059d7b4aeb0\endcsname#1{%
Gospel acclamation: Cf. John 15:16 & Narrated event: The Last Supper & Disputed: \chronologyannotationreach{John.15.16}{Matt.26.20-29 Mark.14.17-25 Luke.22.14-38 John.13.1- John.14 John.15 John.16 John.17.1-26}{true}\chronologyannotationclaim{life-of-christ.last-supper}{narrated-event}{catholic-traditional-v1}{disputed}{the fourteenth day of Nisan}{the fourteenth day of Nisan}\\*[.1em]%
}
\expandafter\def\csname triptychchronologyorientation@row-885908f6c6cd158f\endcsname#1{%
 & Composition: John; Ephesus traditionally & Disputed: \chronologyannotationreach{John.15.16}{John}{true}\chronologyannotationclaim{composition.gospel-of-john}{composition}{catholic-traditional-v1}{disputed}{from the year 90 to 100 (approximately)}{c.~A.D.~90--100}; \chronologyannotationreach{John.15.16}{John}{true}\chronologyannotationclaim{composition.gospel-of-john}{composition}{catholic-traditional-v1}{disputed}{96 or one of the succeeding years}{A.D.~96--100}\\[.3em]%
}
\expandafter\def\csname triptychchronologyorientation@row-9d0a28314a61693e\endcsname#1{%
Communion antiphon: Corinthians option: Cf. 1 Corinthians 10:17, 1 Cor. 10:16 & Composition: The First Epistle to the Corinthians & Disputed: \chronologyannotationreach{1Cor.10.16}{1Cor}{true}\chronologyannotationreach{1Cor.10.17}{1Cor}{true}\chronologyannotationclaim{composition.first-epistle-to-the-corinthians}{composition}{catholic-traditional-v1}{disputed}{(1 and 2 Corinthians; Galatians), 56}{A.D.~56}; \chronologyannotationreach{1Cor.10.16}{1Cor}{true}\chronologyannotationreach{1Cor.10.17}{1Cor}{true}\chronologyannotationclaim{composition.first-epistle-to-the-corinthians}{composition}{catholic-traditional-v1}{disputed}{about Easter A.D. 57}{A.D.~57}\\[.3em]%
}
\expandafter\def\csname triptychchronologyorientation@row-7bc2a7d9295607c3\endcsname#1{%
Second reading: Phil. 4:6--9 & Composition: Philippians; Rome to Philippi in the cited account & Disputed: \chronologyannotationreach{Phil.4.6}{Phil}{true}\chronologyannotationreach{Phil.4.7}{Phil}{true}\chronologyannotationreach{Phil.4.8}{Phil}{true}\chronologyannotationreach{Phil.4.9}{Phil}{true}\chronologyannotationclaim{composition.epistle-to-the-philippians}{composition}{catholic-traditional-v1}{disputed}{(Philemon; Colossians; Ephesians; Philippians), 61}{A.D.~61}; \chronologyannotationreach{Phil.4.6}{Phil}{true}\chronologyannotationreach{Phil.4.7}{Phil}{true}\chronologyannotationreach{Phil.4.8}{Phil}{true}\chronologyannotationreach{Phil.4.9}{Phil}{true}\chronologyannotationclaim{composition.epistle-to-the-philippians}{composition}{catholic-traditional-v1}{disputed}{at Rome (A.D. 62-64)}{A.D.~62--64}\\[.3em]%
}

% Generated. Do not edit.
% Complete orientation-v3 table; wording, dates and order belong to the corpus.
\begin{chronologysharedtable}
\chronologyorientationrow{row-01242fba70f5222d}{}
\midrule
\chronologyorientationrow{row-2a5f76fbcf6ad0bd}{}
\chronologyorientationrow{row-af9138e4ceea5e79}{}
\midrule
\chronologyorientationrow{row-34474a8e5d3d77f6}{}
\midrule
\chronologyorientationrow{row-58ef6f1ff512a823}{}
\midrule
\chronologyorientationrow{row-8dc930e93054888b}{}
\midrule
\chronologyorientationrow{row-a6eb6120b16899b0}{}
\chronologyorientationrow{row-0bfa70ac74041a8c}{}
\midrule
\chronologyorientationrow{row-032e4059d7b4aeb0}{}
\chronologyorientationrow{row-885908f6c6cd158f}{}
\midrule
\chronologyorientationrow{row-9d0a28314a61693e}{}
\midrule
\chronologyorientationrow{row-7bc2a7d9295607c3}{}
\end{chronologysharedtable}
{\footnotesize Psalm numbering: Vg. = Vulgate; Heb. = Hebrew.\par}


\subsection{Evidence for each date}
\label{sec:date-location-evidence}
{\small For each passage, as the Scriptural Date and Location table names it: what is dated, how the corpus weighs each answer, the source's own words, and the sources cited. A subject the corpus does not date is reported as such and given no date.\par}
{\small
\par\smallskip\noindent\textit{Entrance antiphon: Esth. 13:9--11}\par\nobreak
\begin{itemize}[leftmargin=1.2em,itemsep=1pt,topsep=1pt,parsep=0pt]
\item \textbf{Composition} (preferred). The Book of Esther: B.C.~465--425 (Conjectural early dating; later date possible) (``at the end of the reign of Xerxes I (485-465 B.C.) or during the reign of his son Artaxerxes I (465-425 B.C.)''; The Catholic Encyclopedia, Volume 5 (New York, 1909)). Place: story at Susa; writing in Persia conjectured (The Catholic Encyclopedia, Volume 5 (New York, 1909); The Holy Bible, Douay-Rheims (Challoner revision), Project Gutenberg eBook 1581: Esth. 1:2).
\end{itemize}
\par\smallskip\noindent\textit{Responsorial psalm: Ps. 79:9,12--16,19--20 (Vg.; Heb. 80)}\par\nobreak
\begin{itemize}[leftmargin=1.2em,itemsep=1pt,topsep=1pt,parsep=0pt]
\item \textbf{Historical setting}. No date: the chronology corpus has not yet dated this (research pending).
\item \textbf{Composition} (preferred). Psalms: Before c.~165 B.C. (``before the Maccabean period, around 165 B.C.; no individual psalm can be dated securely''; Catholic critical chronology; New American Bible, Revised Edition --- USCCB web state (2026), Psalms introduction).
\end{itemize}
\par\smallskip\noindent\textit{First reading: Is. 5:1--7}\par\nobreak
\begin{itemize}[leftmargin=1.2em,itemsep=1pt,topsep=1pt,parsep=0pt]
\item \textbf{Traditional attribution} (preferred). Isaias (ministry in Souvay's traditional account): B.C.~740--701 (``from the closing year of Ozias, King of Juda (740 B.C.); no prophecies appear to be later than 701 B.C.''; The Catholic Encyclopedia, Volume 8 (New York, 1910)). Place: Judah and Jerusalem (The Holy Bible, Douay-Rheims (Challoner revision), Project Gutenberg eBook 1581: Is. 1:1).
\end{itemize}
\par\smallskip\noindent\textit{Responsorial psalm: Is. 5:7}\par\nobreak
\begin{itemize}[leftmargin=1.2em,itemsep=1pt,topsep=1pt,parsep=0pt]
\item \textbf{Traditional attribution}. As under First reading: Is. 5:1--7.
\end{itemize}
\par\smallskip\noindent\textit{Communion antiphon: Lamentations option: Lam. 3:25}\par\nobreak
\begin{itemize}[leftmargin=1.2em,itemsep=1pt,topsep=1pt,parsep=0pt]
\item \textbf{Composition} (preferred). The poems of Lamentations after Jerusalem's destruction: Responses to Jerusalem's destruction in 587 B.C. (``after Jerusalem's destruction in 587 B.C.; grouped later at an unknown date''; Catholic critical chronology; New American Bible, Revised Edition --- USCCB web state (2026), Lamentations introduction).
\end{itemize}
\par\smallskip\noindent\textit{Gospel: Matt. 21:33--43}\par\nobreak
\begin{itemize}[leftmargin=1.2em,itemsep=1pt,topsep=1pt,parsep=0pt]
\item \textbf{Narrated event}. No date: the chronology corpus has not yet dated this (research pending).
\item \textbf{Composition} (disputed). The Gospel of St Matthew: c.~A.D.~38--45 (``about A.D. 38-45''; The Catholic Encyclopedia, Volume 10 (1911); The Catholic Encyclopedia, Volume 14 (New York, 1912)); c.~A.D.~40--42 (``about the year 40-42''; The Catholic Encyclopedia, Volume 10 (1911)); A.D.~40--45 (``the years 40-45''; The Catholic Encyclopedia, Volume 10 (1911)); c.~A.D.~60--68 (``about the year 60-68''; The Catholic Encyclopedia, Volume 10 (1911)); c.~A.D.~64--67 (``about the years 64-67''; The Catholic Encyclopedia, Volume 10 (1911)); c.~A.D.~50 (``about the year 50''; The Catholic Encyclopedia, Volume 14 (New York, 1912)). Place: Semitic original associated with Palestine (The Catholic Encyclopedia, Volume 10 (1911)).
\end{itemize}
\par\smallskip\noindent\textit{Gospel acclamation: Cf. John 15:16}\par\nobreak
\begin{itemize}[leftmargin=1.2em,itemsep=1pt,topsep=1pt,parsep=0pt]
\item \textbf{Narrated event} (disputed). The Last Supper: the fourteenth day of Nisan (The Catholic Encyclopedia, Volume 8 (New York, 1910)).
\item \textbf{Composition} (disputed). The Gospel of St John: c.~A.D.~90--100 (``from the year 90 to 100 (approximately)''; The Catholic Encyclopedia, Volume 14 (New York, 1912)); A.D.~96--100 (``96 or one of the succeeding years''; The Catholic Encyclopedia, Volume 8 (New York, 1910)). Place: Ephesus traditionally (The Catholic Encyclopedia, Volume 8 (New York, 1910)).
\end{itemize}
\par\smallskip\noindent\textit{Communion antiphon: Corinthians option: Cf. 1 Corinthians 10:17, 1 Cor. 10:16}\par\nobreak
\begin{itemize}[leftmargin=1.2em,itemsep=1pt,topsep=1pt,parsep=0pt]
\item \textbf{Composition} (disputed). The First Epistle to the Corinthians: A.D.~56 (``(1 and 2 Corinthians; Galatians), 56''; The Catholic Encyclopedia, Volume 11 (New York, 1911)); A.D.~57 (``about Easter A.D. 57''; The Catholic Encyclopedia, Volume 4 (New York, 1908)).
\end{itemize}
\par\smallskip\noindent\textit{Second reading: Phil. 4:6--9}\par\nobreak
\begin{itemize}[leftmargin=1.2em,itemsep=1pt,topsep=1pt,parsep=0pt]
\item \textbf{Composition} (disputed). The Epistle to the Philippians: A.D.~61 (``(Philemon; Colossians; Ephesians; Philippians), 61''; The Catholic Encyclopedia, Volume 11 (New York, 1911)); A.D.~62--64 (``at Rome (A.D. 62-64)''; The Catholic Encyclopedia, Volume 12 (1911)). Place: Rome to Philippi in the cited account (The Catholic Encyclopedia, Volume 12 (1911); The Holy Bible, Douay-Rheims (Challoner revision), Project Gutenberg eBook 1581: Phil. 1:1).
\end{itemize}
\par}
